\section{Envelope gaps as probability problems}\label{sec:foundations}
\subsection{The scalar graph and the specified termwise relaxation}
Let $B=\prod_{i=1}^n[\ell_i,u_i]$ be a finite box. Fixed coordinates are
removed when necessary. A function is \emph{multiaffine} if it is affine in
each coordinate separately; equivalently, it is a multilinear polynomial
with no repeated variable in a monomial. For a continuous scalar function
$f$ on $B$, let $\vex_B f$ be its greatest convex underestimator and
$\cav_B f$ its least concave overestimator. The vertical section of the
convex hull of its graph is
\[
 \{y:(x,y)\in\conv\{(z,f(z)):z\in B\}\}
 =[\vex_B f(x),\cav_B f(x)].
\]
Define the \emph{scalar graph-hull gap} and, for a specified decomposition
$f=\sum_{e\in E}f_e$, the \emph{termwise gap} by
\begin{align}
 \hgap_B(f;x)&=\cav_B f(x)-\vex_B f(x),\label{eq:hgap}\\
 \tgap_B(f;x)&=\sum_{e\in E}
 [\cav_B f_e(x)-\vex_B f_e(x)].\label{eq:tgap}
\end{align}
The notation $\tgap_B$ always refers to the displayed decomposition.
It is the width obtained by giving each scalar factor its exact graph hull
and imposing $y=\sum_e y_e$. It can change when a factor is split into
smaller pieces, whereas $\hgap_B$ does not. Always,
\begin{equation}\label{eq:gap-order}
 0\le\hgap_B(f;x)\le\tgap_B(f;x),
\end{equation}
because the sum of convex underestimators is a convex underestimator of
the sum, and similarly for concave overestimators. Ratios are taken only
where $\hgap_B(f;x)>0$; gap inequalities apply at every point.

For a positive multilinear polynomial, the factors are its displayed
monomials $a_e\prod_{i\in e}x_i$, with $a_e\ge0$. For a signed bilinear
function, they are the signed products $a_{ij}x_ix_j$. A bilinear factor
has the familiar four-inequality McCormick hull. For a higher-degree factor,
its exact hull and a recursive McCormick relaxation are different objects.
A recursive product relaxation is exact for one monomial on a unit box,
but need not be exact on a general positive box
\citep[Sections 2--3]{Luedtke2012}. A bound for \eqref{eq:tgap} on that
box therefore does not automatically bound the gap of a recursive bilinear
factorization.

Adding an affine function to a factor, or adding a separate affine factor,
changes neither gap. An affine bijection between boxes preserves both
envelopes after composition and hence preserves the gaps for the transformed
factors. These facts follow from the graph-hull representation.

\subsection{Vertex distributions and compatible marginals}
Let $\mathcal V(B)$ denote the vertices of $B$, and let $\laws_B(x)$ be
the distributions on $\mathcal V(B)$ with mean $x$. On the unit cube these
are joint Bernoulli distributions with the prescribed singleton means;
coordinates need not be independent.

\begin{lemma}[Vertex representation]\label{lem:vertex-law}
For a multiaffine $f$ on $B$,
\begin{align}
 \vex_B f(x)&=\min_{P\in\laws_B(x)}\E_P f(X),\label{eq:vex-law}\\
 \cav_B f(x)&=\max_{P\in\laws_B(x)}\E_P f(X).\label{eq:cav-law}
\end{align}
Both extrema are attained. The envelopes are continuous piecewise-affine
functions on $B$.
\end{lemma}
\begin{proof}
Write $x_i=\ell_i+(u_i-\ell_i)p_i$. Independently choose $X_i=u_i$ with
probability $p_i$ and $X_i=\ell_i$ otherwise. Multiaffinity gives
$\E X=x$ and $\E f(X)=f(x)$. Every graph point therefore belongs to the
convex hull of vertex graph points, and the reverse inclusion in the full
graph hull is immediate. The vertical sections give the formulas. The
feasible distributions form a nonempty compact polytope, so both extrema
are attained. The lower and upper surfaces of the vertex graph polytope
are respectively maxima and minima of finitely many affine functions on
$B$, which also proves continuity at its boundary.
\end{proof}

These representations are classical \citep[Section 3]{Luedtke2012}.
They expose the compatibility issue: different factors may use different
minimizing distributions in \eqref{eq:tgap}, whereas \eqref{eq:vex-law}
requires one distribution for their sum.

For $m_e(x)=\prod_{i\in e}x_i$ with $e\ne\varnothing$, the unit-cube
formulas are
\begin{equation}\label{eq:monomial-envelopes}
 \cav m_e(x)=\min_{i\in e}x_i,\qquad
 \vex m_e(x)=\max\{0,\sum_{i\in e}x_i-|e|+1\}.
\end{equation}
To prove them, interpret $\E m_e(X)$ as the probability of joint success.
The upper bound is attained by $X_i=\ind\{U\le x_i\}$ for one uniform
$U\in[0,1]$. For the lower bound, place failure intervals of lengths
$1-x_i$ consecutively on the unit circle. If their total length is below
one they are disjoint, and if it is at least one they cover the circle.
Their union therefore has measure $\min\{1,\sum_{i\in e}(1-x_i)\}$.
Each interval has its prescribed measure, so the complement gives the
claimed minimum joint-success probability. Unused coordinates can be
sampled independently to complete the mean vector.

\begin{proposition}[Common upper attainment and deficiencies]
\label{prop:deficiency}
Let $f(x)=\sum_e a_e m_e(x)$ on $[0,1]^n$, with $a_e\ge0$. Discard
supports of size at most one. For each remaining support choose an anchor
$i_e$ of smallest mean and put
\[
 u_e=x_{i_e},\quad S_e=\sum_{j\in e\setminus\{i_e\}}(1-x_j),
 \quad t_e=\min\{u_e,S_e\},\quad
 D_e(X)=X_{i_e}-\prod_{i\in e}X_i.
\]
Then $D_e$ is the nonnegative indicator that the anchor succeeds and some
other support coordinate fails. Moreover,
\begin{align}
 \cav f(x)&=\sum_e a_eu_e,\label{eq:common-upper}\\
 \tgap(f;x)&=\sum_e a_et_e,\label{eq:term-deficiency}\\
 \hgap(f;x)&=\max_{P\in\laws_{[0,1]^n}(x)}
                     \sum_e a_e\E_P D_e(X).\label{eq:hull-deficiency}
\end{align}
\end{proposition}
\begin{proof}
The common-threshold law attains every single-product upper envelope
simultaneously. Positivity gives \eqref{eq:common-upper}, and subtracting
the two single-product envelopes gives \eqref{eq:term-deficiency}.
Finally, every admissible law has $\E X_{i_e}=u_e$, so minimizing the
expected polynomial value is equivalent to maximizing the weighted
sum of deficiencies.
\end{proof}

The common-upper statement is classical
\citep[Theorems 4--5 in the authors' manuscript]{Luedtke2012}. A distribution
used to prove a gap bound must work for all terms together; independently
optimizing the distribution for each term merely recovers the termwise
relaxation. Mixtures of admissible laws preserve the means, and nonnegative
deficiencies allow the guarantees of different laws to be added.

\begin{lemma}[Two classes handled by independence]\label{lem:easy-terms}
Put $c_0=1-e^{-1}$. Independent Bernoulli rounding gives
$\E D_e\ge c_0t_e/2$ if either every support mean exceeds $1/2$, or at
least two support means are at most $1/2$.
\end{lemma}
\begin{proof}
Independence gives $\E D_e=u_e(1-\prod_{j\ne i_e}x_j)$. In the second
case a nonanchor mean is at most $1/2$, giving at least $u_e/2\ge t_e/2$.
In the first case the expression is at least
$u_e(1-e^{-S_e})\ge c_0u_e\min(1,S_e)$. If $S_e\le1$, this is at least
$c_0S_e/2\ge c_0t_e/2$; otherwise it is at least $c_0u_e=c_0t_e$.
\end{proof}

\begin{proposition}[Nonnegative-box transfer]\label{prop:box-transfer}
For a positive multilinear polynomial on $B$ with $0\le\ell_i<u_i$,
substitute $x_i=\ell_i+(u_i-\ell_i)p_i$ and expand the original monomials.
The unit-cube polynomial $\widetilde f$ has nonnegative coefficients and
no larger degree. With its expanded monomials as factors,
\begin{equation}\label{eq:expanded-gap}
 \hgap_B(f;x)=\hgap(\widetilde f;p),\qquad
 \tgap_B(f;x)\le\tgap(\widetilde f;p).
\end{equation}
One common-threshold endpoint law attains the upper envelope of every
original monomial and their sum.
\end{proposition}
\begin{proof}
Expansion coefficients are products of nonnegative endpoints and interval
widths. The full-envelope equality follows from the affine bijection.
For an original factor expanded as $\sum_q g_q$, the sum of the lower
convex envelopes is a convex underestimator and the sum of the upper
concave envelopes is a concave overestimator. Subtract and sum over the
original factors to obtain \eqref{eq:expanded-gap}. Common-threshold
rounding attains all expanded upper envelopes together and therefore also
those of each original factor and their sum.
\end{proof}

Expansion is a proof device and need not be an efficient representation.
It changes the factor incidence structure. Degree bounds transfer through
\eqref{eq:expanded-gap}; bounds based on the original incidence graph
require separate arguments. For bilinear functions, arbitrary box endpoints
cause no such difficulty: an affine change introduces only scaled bilinear
terms and affine summands, regardless of signs.

\subsection{What a spatial certificate counts}\label{subsec:certificate-model}
Let $\mathcal F$ be a nonempty compact feasible set, let $f$ be continuous,
and put $f^*=\min_{x\in\mathcal F}f(x)$. A region oracle assigns a valid
lower bound $\operatorname{LB}(S)$ to feasible points in each region $S$.
A \emph{certificate at target $T_*$} is a finite cover $S_1,\ldots,S_N$
of $\mathcal F$ with $\operatorname{LB}(S_j)\ge T_*$ for all $j$.
Together with a feasible value $U$ and an absolute tolerance $\epsilon\ge0$,
the target $U-\epsilon$ certifies absolute gap at most $\epsilon$.
For $U>0$ and a relative tolerance $0\le\theta<1$, the target
$(1-\theta)U$ certifies $(U-f^*)/U\le\theta$. A size lower bound may
grant the best incumbent $U=f^*$, requiring $f^*>0$ in this relative
convention. Under these tolerance ranges, either target is nondecreasing
in $U$, so worse incumbents require weakly higher targets.

A complete spatial split tree supplies a cover by its final regions.
A lower bound for arbitrary covers is thus independent of branching order
and node selection within the specified region and oracle model. If a
procedure discards feasible regions by propagation or additional cuts,
those deductions must also be included in the certificate model. A lower
bound for one oracle does not automatically survive stronger coupled cuts,
different branch coordinates, or analytic solution of separate components.

\section{Signed bilinear gaps and interaction density}\label{sec:bilinear}
Let $G=(V,E)$ be a finite simple graph and
$b_a(x)=\sum_{ij\in E}a_{ij}x_ix_j$ on $[0,1]^V$, with arbitrary real
coefficients. Zero-coefficient edges may be removed. Let $c^*(a)$ be the
least $c\ge0$ with $\tgap(b_a;x)\le c\hgap(b_a;x)$ everywhere, and set
$c^*(0)=0$. Define
\begin{equation}\label{eq:density-definition}
 \dens(G)=\max_{\varnothing\ne W\subseteq V}\frac{|E(G[W])|}{|W|},
\end{equation}
with value zero for an empty graph. This is half the maximum subgraph
average degree, not fractional arboricity, whose denominator is $|W|-1$.
Let $d_G$ be degeneracy and $\Delta_G$ maximum degree. Then
\begin{equation}\label{eq:density-comparisons}
 d_G/2\le\dens(G)\le d_G,\qquad \dens(G)\le\Delta_G/2.
\end{equation}
The first lower bound uses a subgraph of minimum degree $d_G$; the upper
bounds use a degeneracy orientation and the degree sum identity.

\begin{theorem}[Bilinear upper bounds]\label{thm:bilinear-upper}
For every real coefficient vector on $G$,
\begin{equation}\label{eq:bilinear-upper}
 c^*(a)\le\min\{4\sqrt{\dens(G)},\ 4\sqrt{d_G},\ 2\sqrt{\Delta_G}\}.
\end{equation}
If $G$ is bipartite with parts $P,Q$ and has an edge, let $\Delta_P$ and
$\Delta_Q$ be the maximum degrees in the respective parts. Then
\begin{equation}\label{eq:bipartite-upper}
 c^*(a)\le\sqrt{2\min\{\Delta_P,\Delta_Q\}}
       \le\sqrt{2\min\{|P|,|Q|\}}.
\end{equation}
The factor $\sqrt2$ in the first bipartite bound is sharp. The conclusions
transfer to arbitrary finite boxes after removing fixed coordinates.
\end{theorem}
The proof refines the cut argument of \citet{Boland2017}. The density order
also follows from older operator theory, as shown in
Section~\ref{subsec:schur}; no new general graph-norm principle is claimed.

\subsection{An exact characterization by induced cuts}
For $W\subseteq V$, set
\[
 Q_{a,W}(s)=\sum_{ij\in E(G[W])}a_{ij}s_is_j,\qquad
 L_W(a)=\sum_{ij\in E(G[W])}|a_{ij}|,
\]
and define
\begin{equation}\label{eq:cut-range}
 R_W(a)=\tfrac12\left(\max_{s\in\{-1,1\}^W}Q_{a,W}(s)
                         -\min_{s\in\{-1,1\}^W}Q_{a,W}(s)\right).
\end{equation}
A cut encoded by $s$ has weight $\tfrac12\sum_{ij}a_{ij}-\tfrac12Q_{a,W}(s)$,
so $R_W$ is exactly the maximum cut weight minus the minimum cut weight.
If $L_W>0$ then $R_W>0$: distinct degree-two sign characters are orthogonal
under uniform signs, and their nonzero combination cannot be constant.
Set $L_W/R_W=0$ when both vanish.

\begin{lemma}[Induced-cut characterization]\label{lem:induced-cut}
For every coefficient vector,
\begin{equation}\label{eq:cut-characterization}
 c^*(a)=\max_{W\subseteq V}\frac{L_W(a)}{R_W(a)}.
\end{equation}
Each nonzero induced ratio is attained at $x_i=1/2$ on $W$ and $x_i=0$
outside $W$.
\end{lemma}
\begin{proof}
At that point any admissible vertex law fixes outside coordinates to zero.
Write $X_i=(1+s_i)/2$ on $W$. Since $\E s_i=0$,
\[
 \E b_a(X)=\tfrac14\left(\sum_{ij\in E(G[W])}a_{ij}+\E Q_{a,W}(s)\right).
\]
An extremizing sign vector and its negative, each of probability $1/2$,
attain either extreme of $Q_{a,W}$ with the prescribed means. Thus
$\hgap=R_W/2$, whereas each active product has width $1/2$ and
$\tgap=L_W/2$. These identities hold at every point of
$\{0,1/2,1\}^V$, with $W$ its set of half-valued coordinates: fixing
coordinates to one instead of zero adds only affine terms on $W$.

To extend the comparison, the width of the product $x_ix_j$ is
\[
 \min(x_i,x_j)-\max(0,x_i+x_j-1)
 =\min(x_i,x_j,1-x_i,1-x_j).
\]
Partition the cube by the hyperplanes $x_i=x_j$, $x_i+x_j=1$ ($i\ne j$),
and $x_i=1/2$. The termwise gap is affine on every closed cell.
Every cell vertex is in $\{0,1/2,1\}^V$. Indeed, active equations fix a
coordinate to $0,1/2$, or $1$, or relate two coordinates by equality or
complementation. In a connected component of these relations, one fixed
coordinate determines all values. If none is fixed, a cycle with an odd
number of complementations forces the value $1/2$. Otherwise a free
parameter remains and a small perturbation preserves all active equations,
contradicting extremality. This accounts for every cell vertex.

The graph-hull gap is concave, being $\cav b_a-\vex b_a$. Hence
$\tgap-c\hgap$ is convex on each cell for $c\ge0$. Its value at any
convex combination of cell vertices is at most the combination of its
values there. The grid-point bound therefore proves the upper direction
of \eqref{eq:cut-characterization}, and the attained ratios prove the other.
\end{proof}
The grid identities and upper implication are those of
\citet[Lemma 1 and Corollary 1]{Boland2017}, following
\citet{Luedtke2012}. The cell argument above supplies the boundary details.

\subsection{Cut ranges, row norms, and fractional orientations}
Let $A$ be the symmetric weighted adjacency matrix on $W$, with zero
diagonal. For a rectangular matrix define
$\norm{M}_{\infty\to1}=\max_{u\text{ sign}}\norm{Mu}_1
=\max_{u,v\text{ sign}}v^{\mathsf T}Mu$; empty matrices have norm zero.

\begin{lemma}[Polarization and row estimates]\label{lem:cut-row}
Writing $R=R_W(a)$,
\begin{align}
 R&=\max_{S\subseteq W}\norm{A_{S,W\setminus S}}_{\infty\to1},
    \label{eq:cut-polarization}\\
 R&\ge\frac1{\sqrt2}\max_{S\subseteq W}
          \sum_{i\in S}\norm{a_{i,W\setminus S}}_2,
    \label{eq:cut-khinchin}\\
 R&\ge\frac14\sum_{i\in W}\norm{a_i}_2.\label{eq:cut-full-rows}
\end{align}
\end{lemma}
\begin{proof}
For signs $s,t$, put $u=(s+t)/2$, $v=(s-t)/2$. Their supports are
complementary, and symmetry gives $Q(s)-Q(t)=2v^{\mathsf T}Au$.
Conversely, arbitrary signs on complementary supports arise from
$s=u+v$, $t=u-v$. Maximizing half the difference proves
\eqref{eq:cut-polarization}.

The sharp real Khinchin inequality of \citet{Szarek1976} states that
independent uniform signs satisfy
\begin{equation}\label{eq:khinchin}
 \E\left|\sum_j c_j\varepsilon_j\right|
       \ge\frac1{\sqrt2}\left(\sum_j c_j^2\right)^{1/2}.
\end{equation}
Apply it to every row of $A_{S,T}$ for a fixed partition $S,T$, and sum.
The best sign vector performs at least as well as its average, proving
\eqref{eq:cut-khinchin}. Transposition gives the same estimate for side $T$.

Choose a partition locally maximizing the cut weight with nonnegative
weights $a_{ij}^2$. At every vertex at least half its squared incident
weight crosses, since otherwise moving the vertex improves the cut. Thus
\[
 X=\sum_{i\in S}\norm{a_{i,T}}_2,\qquad
 Y=\sum_{i\in T}\norm{a_{i,S}}_2
 \quad\text{satisfy}\quad
 X+Y\ge\frac1{\sqrt2}\sum_i\norm{a_i}_2.
\]
The two row estimates yield
$R\ge\max(X,Y)/\sqrt2\ge(X+Y)/(2\sqrt2)$, which proves
\eqref{eq:cut-full-rows}.
\end{proof}

\begin{lemma}[Fractional orientation]\label{lem:fractional-orientation}
Every finite graph $H$ admits nonnegative fractions
$\theta_{ij}+\theta_{ji}=1$ on each edge such that
$\sum_j\theta_{ij}\le\dens(H)$ for every vertex. This is the least
possible uniform bound on the vertex loads.
\end{lemma}
\begin{proof}
For $m=|E(H)|>0$, form a flow network with a source, edge nodes, vertex
nodes, and a sink. Source-to-edge capacities are one; arcs from an edge
node to its endpoints have capacity $m+1$; vertex-to-sink capacities are
$t$. A cut of capacity at most $m$ cannot cut an edge-to-endpoint arc.
For a fixed set $U$ of vertex nodes on the source side, placing all edge
nodes internal to $U$ on that side minimizes capacity, giving
$m-|E(H[U])|+t|U|$. At $t=\dens(H)$ every cut has capacity at least $m$.
Max-flow/min-cut supplies a flow saturating the source arcs; outgoing flows
at each edge define its fractions. Conversely, the loads over $U$ sum to
at least $|E(H[U])|$, proving minimality. For $m=0$ use zero loads.
\end{proof}

\begin{proof}[Proof of Theorem~\ref{thm:bilinear-upper}]
Fix an induced graph $H$ with nonempty weighted support and abbreviate
$L=\sum_{ij}|a_{ij}|$, $R=R_{V(H)}(a)$. Weighted Cauchy--Schwarz gives
\[
 \sum_j\theta_{ij}|a_{ij}|
 \le\left(\sum_j\theta_{ij}\right)^{1/2}
      \left(\sum_j\theta_{ij}a_{ij}^2\right)^{1/2}
 \le\sqrt{\dens(H)}\norm{a_i}_2.
\]
Summing counts each absolute edge weight once. By
\eqref{eq:cut-full-rows}, $L\le4\sqrt{\dens(H)}R$.
Also $\norm{a_i}_2\ge\norm{a_i}_1/\sqrt{\Delta_H}$ at every nonisolated
vertex, and $\sum_i\norm{a_i}_1=2L$. The same row bound therefore gives
$R\ge L/(2\sqrt{\Delta_H})$. The graph parameters do not increase in
induced subgraphs, so Lemma~\ref{lem:induced-cut} proves
\eqref{eq:bilinear-upper}.

For bipartite $H$ use its inherited parts $P',Q'$ in
\eqref{eq:cut-khinchin}. Every edge crosses, giving
\[
 R\ge\frac1{\sqrt2}\sum_{i\in P'}\norm{a_i}_2
       \ge\frac{L}{\sqrt{2\Delta_P}}.
\]
Transpose to obtain the bound with $\Delta_Q$ and apply the same induced
characterization. On a unit-weight four-cycle with one negative edge,
$L=4$, $R=2$, and $\Delta_P=\Delta_Q=2$. Its center ratio attains two,
so the stated bipartite constant is sharp. An affine box transformation
scales each bilinear coefficient by the two interval widths and introduces
only affine summands, proving the final assertion.
\end{proof}

The degeneracy bound has a distinct direct proof. Orient along a degeneracy
order with outdegree at most $d_G$ and color vertices in reverse order.
Place a vertex opposite the side receiving at least half the squared weight
of its outgoing edges. Its crossing row norm is at least its outgoing row
norm divided by $\sqrt2$. Summing and applying Cauchy--Schwarz to at most
$d_G$ outgoing entries gives $X+Y\ge L/\sqrt{2d_G}$. The final step of
Lemma~\ref{lem:cut-row} yields $R\ge L/(4\sqrt{d_G})$. This proof is
weaker than the density theorem but uses an ordinary acyclic orientation.

\subsection{Sharp order on every interaction graph}
Define
\[
 \Gamma(G)=\sup_{a\in\R^E\setminus\{0\}}c^*(a),\qquad
 \Gamma_\pm(G)=\max_{a\in\{-1,1\}^E}c^*(a),
\]
with both values zero for edgeless graphs.

\begin{theorem}[Worst coefficients on a fixed graph]\label{thm:graph-density}
For every graph with an edge,
\begin{equation}\label{eq:graph-density}
 \max\left\{1,\sqrt{\frac{\dens(G)}{2\log2}}\right\}
 \le\Gamma_\pm(G)\le\Gamma(G)\le4\sqrt{\dens(G)}.
\end{equation}
Requiring every coefficient to be nonzero does not change $\Gamma(G)$.
It also equals
\begin{equation}\label{eq:center-supremum}
 \sup_{a\ne0}L_V(a)/R_V(a).
\end{equation}
The lower bound has a full-support signing and a witness in $\{0,1/2\}^V$.
\end{theorem}
\begin{proof}
On an $h$-vertex, $m$-edge graph with $m>0$, choose independent edge signs
$\sigma_e$. For fixed vertex signs $s$ and $\lambda>0$,
$\E e^{\lambda Q_\sigma(s)}=\cosh(\lambda)^m\le e^{m\lambda^2/2}$,
and similarly for $-\lambda$. Put $Z=\max_s|Q_\sigma(s)|$. Global sign
reversal preserves $Q_\sigma$, so
\[
 e^{\lambda Z}\le\sum_{s:s_1=1}
       (e^{\lambda Q_\sigma(s)}+e^{-\lambda Q_\sigma(s)}).
\]
Taking expectations and using Jensen gives
$\E Z\le h\log2/\lambda+m\lambda/2$. At
$\lambda=\sqrt{2h\log2/m}$, at least one signing satisfies
\begin{equation}\label{eq:random-sign}
 \norm{Q_\sigma}_\infty\le\sqrt{2mh\log2}.
\end{equation}
The values at different sign vectors need not be independent.

Apply this to a densest induced graph $H=G[W]$. Since $L_W=m$ and
$R_W\le\norm{Q_\sigma}_\infty$,
$L_W/R_W\ge\sqrt{\dens(G)/(2\log2)}$. Give all remaining edges arbitrary
signs; they do not affect the witness face of Lemma~\ref{lem:induced-cut}.
The lower bound one follows from $R_W\le L_W$, and the upper bound is
Theorem~\ref{thm:bilinear-upper}.

Call the supremum in \eqref{eq:center-supremum} $C$. The full center shows
$\Gamma\ge C$. Conversely, every induced ratio is a whole-graph ratio
after zeroing all coefficients outside the induced subgraph. Thus
$\Gamma\le C$. On the compact set $L_V(a)=1$, $R_V$ is continuous and
strictly positive, so the ratio attains a maximum. Perturb zero entries
of a maximizer to nonzero values. Continuity of the full-graph ratio shows
that the full-support supremum is also $C$. This does not assume continuity
of $c^*$ as edges disappear or attainment under the full-support restriction.
\end{proof}

Uniform boundedness of all signed factors in a graph family is therefore
equivalent to uniform boundedness of its induced density. No closure
assumption on the family is needed: fixing outside coordinates exposes a
densest subgraph. Global average density is insufficient. Attaching $q^2$
leaves to one vertex of $K_{q,q}$ produces a connected graph with total
edge-to-vertex ratio less than two but induced density at least $q/2$.

The density is unweighted. Scaling all coefficients by a nonzero scalar
preserves the gap ratio, whereas a sum of coefficient magnitudes scales.
A single edge of arbitrarily small magnitude therefore rules out replacing
$\dens$ by $\max_W\sum_{e\in E(W)}|a_e|/|W|$. The row estimates remain
coefficient-sensitive certificates; a uniform gap bound must use them on
every induced subgraph, not only the full graph.

The order may equivalently be expressed using degeneracy or arboricity.
If $\alpha_G$ is the minimum number of forests partitioning the edges,
then $\dens(G)\le\alpha_G\le d_G\le2\dens(G)$. For the middle inequality,
label outgoing edges in a degeneracy orientation with distinct labels
from $1$ to $d_G$. Each label class is a forest: an undirected cycle would
have an earliest vertex with two outgoing edges of that label. The other
inequalities follow by counting forest edges and
\eqref{eq:density-comparisons}.

\subsection{Positive coefficients, exactness, and finite constants}
For positive coefficients a proper coloring with $\chi\ge2$ colors gives
the classical stronger bound
\begin{equation}\label{eq:positive-bilinear}
 c^*(a)\le
 \begin{cases}2-2/\chi,&\chi\text{ even},\\
 2-2/(\chi+1),&\chi\text{ odd}.
 \end{cases}
\end{equation}
Choose uniformly $\lfloor\chi/2\rfloor$ colors for one side of a cut.
Every edge crosses with probability
$p_\chi=2\lfloor\chi/2\rfloor\lceil\chi/2\rceil/[\chi(\chi-1)]$.
Thus the maximum cut is at least $p_\chi L$, and the minimum cut is zero.
The same coloring works on every induced subgraph, so
Lemma~\ref{lem:induced-cut} gives \eqref{eq:positive-bilinear}.
This is the bound of \citet[Theorem 8 in the authors' manuscript]{Luedtke2012}.
In particular, the positive bilinear factor is at most two and is one when
its nonzero support is nonempty and bipartite. For the all-positive $K_n$
with $n\ge2$, the center ratio is
$\binom n2/\lfloor n^2/4\rfloor$, tending to two. Hence the worst-over-signs
quantifier in Theorem~\ref{thm:graph-density} is essential.

Assume the effective support graph, obtained by deleting zero-coefficient
edges, has at least one edge. Then $c^*(a)=1$ if and only if every cycle
of this graph
has an even number of positive and an even number of negative edges
\citep[Theorem 4]{Boland2017}. Boland et al.\ identify this as a consequence
of the convex-envelope criterion in Theorem 3.10 of Misener, Smadbeck,
and Floudas, \emph{Dynamically generated cutting planes for mixed-integer
quadratically constrained quadratic programs and their incorporation into
GloMIQO 2}; our attribution to that predecessor is through Boland et al.,
who give the alternative cut proof used here.
Indeed, $R_V=L_V$ exactly when one cut
contains all positive edges and no negative edges, and another contains
all negative edges and no positive edges. An edge set is a cut exactly
when it intersects every cycle evenly. Necessity follows by traversing a
cycle. For sufficiency, label vertices by the parity of specified edges
on a path from a root in each component; cycle parity makes the label
path-independent and supplies the cut. The criterion passes to induced
subgraphs, proving sufficiency for $c^*=1$. Conversely, $c^*=1$ implies
$R_V=L_V$ by Lemma~\ref{lem:induced-cut}. In particular, every forest with
at least one edge has factor one for every signing. For empty effective
support the gaps coincide at zero and $c^*=0$ by convention.

The best universal constant in $\Gamma(G)\le K_\rho\sqrt{\dens(G)}$
lies in $[2,4]$: the lower bound is the four-cycle with one negative edge,
with density one and factor two. For degeneracy the corresponding interval
is $[\sqrt2,4]$. These are finite-parameter questions, distinct from the
best constants along sequences with density or degeneracy tending to
infinity; the four-cycle does not give a growing-density lower bound.

The bounds imply $c^*\le4\sqrt2$ for series-parallel graphs by degeneracy,
$c^*\le4\sqrt3$ for planar graphs by density, and
$c^*\le\sqrt{2\min(p,q)}$ for $K_{p,q}$. For $K_n$, the maximum-degree
bound is $2\sqrt{n-1}$, compared with $600\sqrt n$ in
\citet[Theorem 2]{Boland2017}. These are scalar vertical-gap comparisons,
not approximation guarantees for a constrained objective.

\subsection{Exact finite full-signing constants}\label{subsec:finite-signings}
For $n\ge2$, distinguish the largest full-center ratio among complete-graph
signings from the largest ratio over all points:
\[
 M_n=\max_{a\in\{-1,1\}^{E(K_n)}}\frac{\binom n2}{R_V(a)},
 \qquad F_n=\Gamma_\pm(K_n).
\]
All edges here have magnitude one. These maxima differ from the
arbitrary-real-coefficient supremum in \eqref{eq:center-supremum}.

\begin{proposition}[Finite exhaustive computation]\label{prop:finite-signings}
The exact values are
\[
\begin{array}{c|rrrrr}
 n&3&4&5&6&7\\\hline
 \min_a R_V(a)&2&4&4&5&8\\
 M_n&3/2&3/2&5/2&3&21/8\\
 F_n&3/2&3/2&5/2&3&3
\end{array}
\]
\end{proposition}
\begin{proof}
The first row is established by the exhaustive integer computation in
Appendix~\ref{app:finite-signings}, which proves completeness of the
enumeration, gives executable code, and supplies attaining signings.
Dividing $\binom n2$ by that row gives $M_n$.
Lemma~\ref{lem:induced-cut} gives
\[
 F_n=\max_{2\le k\le n} M_k.
\]
Indeed, each induced support of a full signing is a full signing of some
$K_k$, giving the upper bound. Conversely, any signing of $K_k$ extends
to a full signing of $K_n$ by assigning signs to all remaining edges,
and its ratio persists on the face fixing the new coordinates to zero.
Since $M_2=1$, this proves the last row.
\end{proof}

In particular $M_7<M_6$ while $F_7=F_6=3$. An explicit witness is the
$K_6$ signing in Table~\ref{tab:finite-witnesses}, extended by a new vertex
with all incident coefficients positive. At the point whose six old
coordinates are $1/2$ and whose new coordinate is zero, its ratio is three;
at the full center its ratio is only $21/10$. The finite computation
optimizes precisely these listed full-signing problems; it makes no claim
about larger $n$ or the optimal constants for arbitrary real coefficients.

\subsection{The relation to Schur-multiplier theory}\label{subsec:schur}
We record the older implication to separate the density order from the
explicit elementary constant. For a real matrix $M$, its Schur multiplier
acts by entrywise multiplication, with norm $\norm{M}_{\mathrm{Schur}}$
measured in the Euclidean operator norm. Define the real projective norm
\[
 \pi(M)=\inf\left\{\sum_k\norm{u_k}_\infty\norm{v_k}_\infty:
                         M=\sum_k u_kv_k^{\mathsf T}\right\}.
\]
The established results of
\citet[Theorems 1.2 and 2.4]{DavidsonDonsig2007} give
\begin{equation}\label{eq:schur-input}
 \norm{M}_{\mathrm{Schur}}\le2\sqrt{\beta(P)},\qquad
 \pi(M)\le K_G\norm{M}_{\mathrm{Schur}},
\end{equation}
for $|M_{ij}|\le1$ supported on a pattern $P$, where $K_G$ is the real
Grothendieck constant and
\[
 \beta(P)=\max_{R,C:\,|R|+|C|>0}
             \frac{|P\cap(R\times C)|}{|R|+|C|}.
\]
The first bound uses the weighted matrix theorem, which permits continuous
row and column bounds. Using the integer pattern theorem instead would
introduce rounding. These are established norm inequalities, not statements
about sizes of convex extended formulations.

For the symmetric adjacency pattern,
\[
 |P\cap(R\times C)|\le|E(G[R\cup C])|+|E(G[R\cap C])|
                   \le\dens(G)(|R|+|C|).
\]
An edge is counted at most once unless both endpoints lie in $R\cap C$,
when both matrix orientations are counted. Taking $R=C$ equal to a densest
set proves $\beta(P)=\dens(G)$. For weighted adjacency $A$, take
$M=\operatorname{sign}(A)$, with zeros where $A$ vanishes. The definition
of $\pi$ gives
\[
 2L_V=\ip{A}{M}\le\pi(M)\norm{A}_{\infty\to1}.
\]
For sign vectors $u,v$, put $p=(u+v)/2$, $q=(u-v)/2$. Symmetry yields
$u^{\mathsf T}Av=2[Q_{a,V}(p)-Q_{a,V}(q)]$. Multiaffinity puts these two
values between the vertex extrema on $[-1,1]^V$, so
$\norm{A}_{\infty\to1}\le4R_V$. Consequently
\begin{equation}\label{eq:old-density-transfer}
 L_V\le4K_G\sqrt{\dens(G)}R_V.
\end{equation}
Apply this on every induced graph to recover the density upper order.
The direct random-sign argument supplies the matching bilinear lower order.

Equivalently, let
$S(G)=\sup_{a\ne0}L_V(a)/\norm{Q_{a,V}}_\infty$ be the real Sidon
constant of the edge-supported degree-two sign characters, with $S(G)=0$
for edgeless graphs. Because $Q$
has mean zero,
$\norm{Q}_\infty/2\le R_V\le\norm{Q}_\infty$ and hence
$S(G)\le\Gamma(G)\le2S(G)$. For bipartite graphs, flipping all signs in
one part negates $Q$, making its range symmetric and giving
$S(G)=\Gamma(G)$. The constant four improves the particular transfer
\eqref{eq:old-density-transfer}; it is not asserted to be the best constant
available in all prior literature.
